\documentclass[11pt,a4paper]{article}
\usepackage[T1]{fontenc}
\usepackage{lmodern}
\usepackage[margin=25mm]{geometry}
\usepackage{amsmath,amssymb}
\usepackage[hidelinks]{hyperref}
\usepackage{microtype}
\pdfinfoomitdate=1
\pdftrailerid{}
\pdfsuppressptexinfo=15
\hypersetup{pdftitle={A ternary counterexample to a proposed mutual-information continuity bound}}
\newcommand{\eps}{\varepsilon}
\newcommand{\TV}{\operatorname{TV}}
\title{A ternary counterexample to a proposed\newline mutual-information continuity bound}
\author{Yongxian Zhang\thanks{School of Computer Science and Engineering,
South China University of Technology, Guangzhou, China.
Email: \href{mailto:mxymmxym1@gmail.com}{mxymmxym1@gmail.com}.
ORCID: \href{https://orcid.org/0009-0000-3864-3536}{0009-0000-3864-3536}.}}
\date{8 October 2026}
\begin{document}
\maketitle
\begin{abstract}
We give explicit classical $3\times3$ probability tables at total variation distance $\eps$ whose mutual-information difference violates the general proposed bound in arXiv:2408.15226v2, Eq.~(106), for every $0<\eps\le1/16$. The same construction defines commuting quantum states. Both marginals change: the version requiring one equal marginal is unresolved by this note. The complete classical counterexample has an independently verified Lean proof; the quantum diagonal embedding and the necessary-leading-coefficient corollary are established by written arguments.
\end{abstract}

\section{Exact target and scope}
Mario Berta, Ludovico Lami and Marco Tomamichel, \emph{Continuity of entropies via integral representations}, propose the general mutual-information continuity inequality
\begin{equation}\label{eq:target}
\left|I(A:B)_\rho-I(A:B)_\sigma\right|
\le h(\eps)+\eps\log\!\left(\min\{d_A^2,d_B^2\}-1\right)
\end{equation}
for nearby bipartite states whose trace distance is at most a sufficiently small $\eps$. This is Eq.~(106) of \href{https://arxiv.org/html/2408.15226v2\#S4}{arXiv:2408.15226v2} and Eq.~(54) of \href{https://arxiv.org/html/2408.15226v1\#S3.SS3}{v1}. The surrounding discussion first considers applications with one equal marginal, then proposes this more general inequality without that restriction.

The construction below refutes the general version with different marginals at arbitrarily small distances. Both marginals change. It does not resolve the version in which one marginal must remain equal.

All logarithms are natural. Changing the logarithm base rescales both sides
and does not alter the strict violation. Write
\begin{equation}
h(\eps)=-\eps\log\eps-(1-\eps)\log(1-\eps),\qquad 0\log0=0.
\end{equation}
For a probability vector $p$, let $H(p)=-\sum_i p_i\log p_i$. For a joint probability table $R$, let $R_A$ and $R_B$ be its row and column marginals and set $I(R)=H(R_A)+H(R_B)-H(R)$.

\section{Construction on a ternary product alphabet}
For every $0<\eps\le1/16$, define probability tables on $\{0,1,2\}\times\{0,1,2\}$ by
\begin{equation}
P_\eps=\begin{pmatrix}
(1-\eps)/2&0&0\\
0&(1-\eps)/2&0\\
0&0&\eps
\end{pmatrix},\qquad
Q_\eps=\begin{pmatrix}
(1-\eps)/2&\eps/2&0\\
\eps/2&(1-\eps)/2&0\\
0&0&0
\end{pmatrix}.
\end{equation}
Every entry is nonnegative, and each table sums to one. The entries at $(0,0)$ and $(1,1)$ agree. The only nonzero entries of $P_\eps-Q_\eps$ are $+\eps$ at $(2,2)$ and $-\eps/2$ at each of $(0,1)$ and $(1,0)$. Hence
\begin{equation}
\TV(P_\eps,Q_\eps)=\frac12\sum_{i,j}|P_\eps(i,j)-Q_\eps(i,j)|
=\frac12(\eps+\eps/2+\eps/2)=\eps.
\end{equation}
The tables also define diagonal density matrices
\begin{equation}
\rho_\eps=\sum_{i,j}P_\eps(i,j)|ij\rangle\langle ij|,\qquad
\sigma_\eps=\sum_{i,j}Q_\eps(i,j)|ij\rangle\langle ij|
\end{equation}
on $\mathbb C^3\otimes\mathbb C^3$. Their quantum trace distance equals this total variation distance, and their von Neumann entropies equal the Shannon entropies of the corresponding diagonal probabilities. Thus the classical commuting construction is also a quantum counterexample.

\section{Exact entropies and mutual information}
Both marginals of $P_\eps$ are
\begin{equation}
(P_\eps)_A=(P_\eps)_B=\big((1-\eps)/2,(1-\eps)/2,\eps\big).
\end{equation}
Its joint distribution is supported on the diagonal. Direct expansion gives
\begin{align}
H((P_\eps)_A)=H((P_\eps)_B)=H(P_\eps)
&=-(1-\eps)\log\frac{1-\eps}{2}-\eps\log\eps\notag\\
&=h(\eps)+(1-\eps)\log2,
\end{align}
so
\begin{equation}
I(P_\eps)=h(\eps)+(1-\eps)\log2.
\end{equation}
Both marginals of $Q_\eps$ are $(1/2,1/2,0)$, so each has entropy $\log2$. Its four nonzero joint entries describe a fair binary input followed by binary crossover probability $\eps$. Equivalently, direct expansion gives
\begin{equation}
H(Q_\eps)=-(1-\eps)\log\frac{1-\eps}{2}-\eps\log\frac{\eps}{2}
=\log2+h(\eps).
\end{equation}
Consequently,
\begin{equation}
I(Q_\eps)=\log2-h(\eps),\qquad
I(P_\eps)-I(Q_\eps)=2h(\eps)-\eps\log2.
\end{equation}

\section{Strict violation at arbitrarily small distances}
Here $d_A=d_B=3$, so the proposed right-hand side is $h(\eps)+\eps\log8$. The signed mutual-information difference minus this bound is
\begin{equation}
I(P_\eps)-I(Q_\eps)-\big(h(\eps)+\eps\log8\big)=h(\eps)-\eps\log16.
\end{equation}
For $0<\eps\le1/16$, monotonicity of the logarithm gives
\begin{equation}
-\eps\log\eps\ge\eps\log16.
\end{equation}
Since $0<1-\eps<1$, also
\begin{equation}
-(1-\eps)\log(1-\eps)>0.
\end{equation}
Adding these proves $h(\eps)-\eps\log16>0$. The signed difference is positive and equals its absolute value. Thus, for \emph{every} $0<\eps\le1/16$,
\begin{equation}
\boxed{\left|I(P_\eps)-I(Q_\eps)\right|>h(\eps)+\eps\log8},\qquad
\TV(P_\eps,Q_\eps)=\eps.
\end{equation}
There is no positive uniform small-distance neighborhood on which the proposed general inequality holds. Given any $\delta>0$, take $\eps=\min\{\delta/2,1/32\}$; then $0<\eps<\delta$, and the construction still violates the bound.

The exact rational choice $\eps=1/32$ gives the strictly positive gap
\begin{equation}
\frac{\log2}{32}-\frac{31}{32}\log\frac{31}{32}>0.
\end{equation}
This positivity is symbolic and does not require numerical certification.

\section{Necessary leading coefficient}
The exact loss $2h(\eps)-\eps\log2$ implies that no modulus of the form $c\,h(\eps)+C\eps$, with finite fixed $C$ and $c<2$, can bound arbitrary ternary mutual-information differences near zero. Indeed,
\begin{equation}
\frac{h(\eps)}{\eps}\ge-\log\eps\longrightarrow+\infty\qquad(\eps\downarrow0),
\end{equation}
so $(2-c)h(\eps)-\eps(\log2+C)>0$ for sufficiently small positive $\eps$. This is a lower bound on the necessary leading coefficient; it does not establish an exact optimal finite-distance modulus.

\section{Formal verification, reproduction and scope}
The complete classical probability-table counterexample is formalized in
Lean 4.34.1. Its public package is fixed to commit
\nolinkurl{363224e0f7cf1d7ea7b4ed2546f8ed372da7b04c}, directory
\nolinkurl{formalizations/mutual-information-continuity-counterexample}.
The definitions are actual finite probability tables, Shannon entropy,
row and column sums, mutual information and half-$L^1$ total variation.
No entropy formula is supplied as an assumed hypothesis.

The two main endpoints in namespace
\nolinkurl{MutualInformationCounterexample} are
\begin{quote}\small
\nolinkurl{counterexample_family}\\
\nolinkurl{counterexamples_arbitrarily_close}.
\end{quote}
The first proves strict violation for every $0<\eps\le1/16$; the second
proves valid witnesses in every positive-distance neighborhood. The
independently compiled \nolinkurl{LiteralSemantics.lean} also checks that
both marginals change.
The fixed verification records report fresh compilation, all 38 owned
declarations and their complete 16,161-declaration dependency closure
replayed in an initially empty kernel at trust level zero. The only logical
axioms are \texttt{propext}, \texttt{Classical.choice} and \texttt{Quot.sound};
an invalid-proof control was rejected. These are prior execution records,
not a claim that this editorial preparation reran Lean.

For reproduction, read the fixed package's \texttt{README.md},
\texttt{case.json}, \texttt{lake-manifest.json} and \texttt{reproduce.py}.
The driver requires the matching compiler and nine pinned dependencies
with their build caches; its explicit command builds a new owned module
and performs the kernel audit. It does not automatically install those
prerequisites. The immutable
\href{https://github.com/mxym/math/releases/tag/mutual-information-continuity-counterexample-v2-lean}{v2 Release}
contains source, full replay records, the compressed dependency graph,
source binding and asset checksums. The earlier
\href{https://github.com/mxym/math/releases/tag/mutual-information-continuity-counterexample-v1}{written v1 Release}
is preserved unchanged.

The quantum diagonal embedding in Section~2 and the necessary coefficient
bound in Section~5 remain written arguments, rather than formalized quantum
spectral-function or asymptotic theorems. The equal-one-marginal question
remains unresolved by this paper. No exhaustive novelty or historical
priority conclusion or external human peer-review certificate is asserted.

\section*{Funding and AI-assisted research methods}
This research was conducted independently and received no external funding.
OpenAI AI assistants assisted mathematical research, proof development,
formalization and internal review. Mathematical claims are supported by
the written proof and the specified formal-verification evidence; AI output
alone is not proof evidence.

\section*{Declaration of generative AI and AI-assisted technologies
in the writing process}
OpenAI AI assistants assisted drafting, organization and LaTeX editing.
The author takes responsibility for the manuscript and its account of the
verification evidence.

\begin{thebibliography}{1}
\bibitem{BLT} Mario Berta, Ludovico Lami and Marco Tomamichel.
\emph{Continuity of entropies via integral representations}.
arXiv:2408.15226. Version~2, Eq.~(106), concluding discussion:
\url{https://arxiv.org/html/2408.15226v2}.
Version~1, Eq.~(54), Section~III.3:
\url{https://arxiv.org/html/2408.15226v1}.
\end{thebibliography}
\end{document}
